\section{Descenso de preesquemas étales: criterios de efectividad}
\label{IX.4}

\ifthenelse{\boolean{orig}}{}
{\enlargethispage{.7cm}}%
\begin{proposition}
\label{IX.4.1}
Sea $g\colon S'\to S$ un morfismo fielmente plano y cuasicompacto.
Entonces~$g$ es un morfismo de descenso efectivo para la categoría
fibrada de los preesquemas étales, separados y de tipo finito sobre otros.
\end{proposition}

En efecto, es un morfismo de descenso para la categoría fibrada en
cuestión, en virtud de~\Ref{IX.3.3} o de~\Ref{VIII.5.2}, a elección.
Queda por demostrar que, si~$X'$ es étale, separado y de tipo finito
sobre~$S'$, y está provisto de un dato de descenso respecto de
$g\colon S'\to S$, este dato es efectivo en la categoría fibrada en
cuestión. Ahora bien, se ve fácilmente que, si~$X$ es un preesquema
sobre~$S$, entonces
\marginpar{236}
es étale sobre~$S$ si y solo si lo es sobre~$S'$ (en virtud de la
Definición~\Ref{IX.1.1} y de \loccit \Ref{VIII.3.6}). Por tanto, es
étale, separado y de tipo finito sobre~$S$ si y solo si~$X'$ lo es
sobre~$S'$; \cf por ejemplo,~\Ref{IX.2.4}. Basta, pues, cerciorarse de
la efectividad del dato de descenso sobre
\ifthenelse{\boolean{orig}}
{$X$}
{$X'$}
para la categoría fibrada de las flechas de~$\Sch$. Ahora bien,~$X'$ es
cuasiafín sobre~$S'$ en virtud de~\Ref{VIII.6.2}
y~\Ref{VIII.6.6}. Se puede concluir entonces utilizando
\Ref{VIII.7.9}. El lector observará, además, que la demostración
exige menos si uno se limita a los preesquemas étales y \emph{finitos}
sobre otros, pues entonces se puede invocar directamente
\Ref{VIII.2.1}.

\begin{corollaire}
\label{IX.4.2}
Sean $g\colon S'\to S$ un morfismo universalmente submersivo, $X'$
\ifthenelse{\boolean{orig}}
{unos $S'$-preesquemas}
{un $S'$-preesquema}
étale, separado y de tipo finito, provisto de un dato de descenso respecto
de~$g$; $S_1\to S$ un morfismo fielmente plano y cuasicompacto; y~$S_1'$
y~$X_1'$ los preesquemas deducidos de~$S'$ y~$X'$ por cambio de base, de
modo que $S_1'\to S_1$ es universalmente submersivo, $X_1'$ es étale,
separado y de tipo finito sobre~$S_1'$, y está provisto de un dato de
descenso respecto de $g_1\colon S_1'\to S_1$. El dato de descenso
sobre~$X'$ es efectivo si y solo si lo es el dato de descenso sobre~$X_1'$.
\end{corollaire}

Esto resulta de la teoría del descenso en las categorías~\cite{IX.D},
teniendo en cuenta~\Ref{IX.4.1} y~\Ref{IX.3.3}.

De manera análoga se demuestra:

\begin{corollaire}
\label{IX.4.3}
Sean $g\colon S'\to S$ un morfismo universalmente submersivo, $X'$ un
$S'$-preesquema étale provisto de un dato de descenso respecto de~$g$, y
$(S_i)$ un recubrimiento de~$S$ por abiertos. El dato de descenso es
efectivo si y solo si, para todo~$i$, lo es el dato de descenso
correspondiente sobre
\ifthenelse{\boolean{orig}}
{$X_i'=X\times_S S_i$}
{$X_i'=X'\times_S S_i$},
respecto del morfismo
$g_i\colon S_i'=S'\times_S S_i\to S_i$.
\end{corollaire}

Este último resultado conduce a aislar un criterio local de efectividad:

\enlargethispage{.3cm}
\begin{proposition}
\label{IX.4.4}
Sea $g\colon S'\to S$ un morfismo de presentación finita en el sentido
de~\Ref{VIII.3.6},
\marginpar{237}
y universalmente submersivo; sea $X'$ un preesquema étale y de
presentación finita sobre~$S'$, provisto de un dato de descenso respecto
de~$g$; y sea, finalmente,~$a$ un punto de~$S$. Para que exista un entorno
abierto~$U$ de~$a$ tal que el dato de descenso correspondiente
sobre~$X_U'=X'\times_S U$ respecto del morfismo
$$g_U\colon S_U'=S'\times_S U\to S_U=U$$
sea efectivo, es necesario y suficiente que el dato de descenso
correspondiente sobre $X_a'=X'\times_S\Spec(\cal{O}_a)$ respecto del
morfismo
$$g_a\colon S_a'=S'\times_S\Spec(\cal{O}_a)\to
S_a=\Spec(\cal{O}_a),$$
sea efectivo.
\end{proposition}

La necesidad es trivial; demostremos la suficiencia. Se dispone, pues, de
un preesquema étale de tipo finito~$X_a$ sobre~$S_a$, y de un isomorfismo
\begin{equation*}
\label{eq:IX.4.*}
\tag{$*$} {X_a'\isomto X_a\times_{S_a}S_a'}
\end{equation*}
compatible con los datos de descenso. De acuerdo con un argumento general
elemental sobre los preesquemas definidos sobre un límite inductivo de
anillos (aquí, los anillos~$A_f$, donde~$A$ es el anillo de un entorno
abierto afín de~$a$ y~$f$ recorre los elementos de~$A$ que no pertenecen
al ideal primo correspondiente a~$a$), se pueden encontrar un entorno
abierto~$U$ de~$a$, un preesquema étale de tipo finito~$X_U$ sobre
$U=S_U$, y un $S_a$-isomorfismo
$X_a\isomto X_U\times_{S_U}S_a$. Además, tomando~$U$ suficientemente
pequeño, se puede suponer que el isomorfismo~\eqref{eq:IX.4.*} proviene de
un isomorfismo
$$
X_U'\isomto X_U\times_{S_U} S_U';
$$
Este último podría no ser compatible con los datos de descenso; sin
embargo, tras restringir~$U$, será compatible con ellos. Esto concluye la
demostración.

\begin{corollaire}
\label{IX.4.5}
Bajo las condiciones de~\Ref{IX.4.4}, el dato de descenso sobre~$X'$ es
efectivo si y solo si, para todo $a\in S$, lo es el dato de descenso
correspondiente sobre~$X_a'$ respecto del morfismo
$S_a'=S'\times_S\Spec(\cal{O}_a)\to\Spec(\cal{O}_a)$. Cuando~$S$ es
localmente noetheriano y~$X'$ está separado sobre~$S'$, en el criterio
precedente se puede sustituir también~$\cal{O}_a$ por su completación.
\end{corollaire}

La
\marginpar{238}
primera afirmación resulta de~\Ref{IX.4.4} y~\Ref{IX.4.3}; la segunda es
entonces consecuencia de~\Ref{IX.4.2}. Utilizando de nuevo~\Ref{IX.4.2}
y el hecho de que, para todo anillo local noetheriano~$A$, se puede
encontrar un anillo local noetheriano completo~$B$ y un homomorfismo local
$A\to B$ tales que~$B$ sea plano sobre~$A$ y que $B/\goth{m}B$ sea una
extensión dada del cuerpo residual $k=A/\goth{m}$ de~$A$, se obtiene:

\begin{corollaire}
\label{IX.4.6}
Bajo las condiciones de~\Ref{IX.4.4}, supongamos además que~$X'$ está
separado sobre~$S'$ y que~$S$ es localmente noetheriano. El dato de
descenso sobre~$X'$ es efectivo si y solo si, para todo preesquema~$S_1$
sobre~$S$ que sea espectro de un anillo local completo con cuerpo residual
algebraicamente cerrado, lo es el dato de descenso correspondiente sobre
$X_1'=X'\times_S S_1$ respecto del morfismo
$g_1\colon S_1'\to S_1$.
\end{corollaire}

\begin{theoreme}
\label{IX.4.7}
Sea $g\colon S'\to S$ un morfismo finito, sobreyectivo y de presentación
finita (esta última hipótesis es consecuencia de las demás si~$S$ es
localmente noetheriano)\footnote{Se puede demostrar que, de hecho, basta
que~$g$ sea un morfismo \emph{entero}, reduciéndose al caso del texto
mediante un procedimiento de paso al límite al estilo de EGA IV~8.}.
Entonces~$g$ es un morfismo de descenso efectivo para la categoría fibrada
de los preesquemas étales, separados y de tipo finito sobre otros.
\end{theoreme}

Hay que demostrar que, si~$X'$ es étale, separado y de tipo finito
sobre~$S'$, y está provisto de un dato de descenso respecto de~$g$, este
dato es efectivo. Utilizando~\Ref{IX.4.3}, se reduce fácilmente al caso en
que~$S$ es noetheriano. Gracias a~\Ref{IX.4.5}, se puede suponer, por
tanto, que~$S$ es el espectro de un anillo local noetheriano y, con mayor
razón, que
$$
\dim S=n<+\infty.
$$
Se razona entonces por inducción sobre~$\dim S=n$; la afirmación es trivial
para $n<0$. Supongamos, pues, que $n\geq 0$ y que el teorema está
demostrado para las dimensiones $n'<n$. En virtud de~\Ref{IX.4.6}, se
reduce al caso en que~$S$ es el espectro de un anillo local completo; por
tanto,~$S'$ es una unión finita de espectros de anillos locales completos.
Se tiene entonces
$$
X'=X_1'\sqcup X_2'
$$
donde~$X_1'$ es \emph{finito} sobre~$S'$ y~$X_2'$ no tiene ningún punto
sobre ninguno de los puntos cerrados
\marginpar{239}
de~$S'$. Consideremos los morfismos
$$
\xymatrix@C=.5cm{ q_1,q_2\colon X'' \ar@/^2pt/[r] \ar@/_2pt/[r] & X' }
$$
correspondientes al dato de descenso, compatibles con
$\xymatrix@C=.5cm{p_1,p_2\colon S'' \ar@/^2pt/[r] \ar@/_2pt/[r] & S'}$.
Se ve inmediatamente que
$$
X''=q_i^{-1}(X_1')\sqcup q_i^{-1}(X_2')\quad i=1,2
$$
es la descomposición canónica análoga de~$X''$ sobre~$S''$, lo que implica
$q_1^{-1}(X_1')=\allowbreak q_2^{-1}(X_1')$; por consiguiente,~$X_1'$ y
$X_2'$ están provistos de datos de descenso inducidos. Sea ahora~$T$ el
abierto de~$S$ complementario de su punto cerrado; entonces
$T'=S'\times_S T$ es la parte de~$S'$ complementaria del conjunto de
puntos cerrados, y~$X_2'$, que se encuentra íntegramente sobre~$T'$, está
provisto de un dato de descenso respecto del morfismo $T'\to T$ inducido
por~$g$. Como este último es finito y sobreyectivo, y
$\dim T<\dim S=n$, este dato de descenso es efectivo por la hipótesis de
inducción. Se ve, pues, que basta demostrar que el dato de descenso
sobre~$X_1'$ es efectivo; por tanto, se puede suponer ahora que~$X'$ es
étale y \emph{finito} sobre~$S'$. (N.B.: el razonamiento por inducción es
innecesario si en el enunciado~\Ref{IX.4.7} se limita uno a los
recubrimientos étales). Sea entonces~$S_0$ el espectro del cuerpo residual
de~$A$, sea $S_0'=S'\times_S S_0$, y definamos análogamente~$S_0''$ y
$S_0'''$ a partir del cuadrado y el cubo fibrados~$S''$ y~$S'''$ de~$S'$
sobre~$S$. En virtud de~\Ref{IX.1.8}, los morfismos $S_0\to S$,
$S_0'\to S'$, etc., inducen equivalencias para las categorías de los
recubrimientos étales de~$S$ y~$S_0$, por una parte, y de~$S'$ y~$S_0'$,
por otra, \ifthenelse{\boolean{orig}}{etc...}{etc.} De los argumentos
formales de la teoría del descenso en las categorías~\cite{IX.D} se deduce
que $g\colon S'\to S$ es un morfismo de descenso efectivo para la
categoría fibrada de los recubrimientos étales si y solo si lo es
$g_0\colon S_0'\to S_0$. Pero este es, en efecto, el caso, por ejemplo
como caso particular de~\Ref{IX.4.1}. Esto concluye la demostración.

\begin{corollaire}
\label{IX.4.8}
La conclusión de~\Ref{IX.4.7} sigue siendo válida si solo se supone que
$S'\to S$ es universalmente submersivo, de tipo finito y cuasifinito,
siempre que~$S$ sea localmente noetheriano.
\end{corollaire}

En
\marginpar{240}
efecto, en virtud de~\Ref{IX.4.6}, se puede suponer que~$S$ es el espectro
de un anillo local noetheriano completo. Entonces, en virtud
de~\Ref{IX.2.5}, existen un morfismo finito y sobreyectivo $S_1\to S$ y un
$S$-morfismo $S_1\to S'$. Como $S_1\to S$ es un morfismo universal de
descenso estricto para la categoría fibrada considerada, en virtud
de~\Ref{IX.4.7}, y $S'\to S$ es un morfismo universal de descenso para
dicha categoría, \Ref{IX.4.8} resulta de los argumentos formales
generales~\cite{IX.D}.

\begin{corollaire}
\label{IX.4.9}
Sea $g\colon S'\to S$ un morfismo de tipo finito, sobreyectivo y
universalmente abierto, con~$S$ localmente noetheriano. Entonces~$g$ es un
morfismo de descenso efectivo para la categoría fibrada de los
preesquemas étales, separados y de tipo finito sobre otros.
\end{corollaire}

Procediendo como en~\Ref{IX.4.7}, se reduce al caso en que~$S$ es el
espectro de un anillo local noetheriano completo~$A$. Sea~$A_1$ una
$A$-álgebra finita, de espectro~$S_1$, tal que $S_1\to S$ sea finito y
\emph{sobreyectivo}, y por tanto un morfismo universal de descenso efectivo
para la categoría fibrada considerada, gracias a~\Ref{IX.4.7}. Se deduce
entonces de los teoremas generales~\cite{IX.D} que~$g$ es un morfismo de
descenso efectivo para la categoría fibrada considerada si y solo si lo es
el morfismo correspondiente
$g_1\colon S_1'=S'\times_S S_1\to S_1$. Como este último satisface las
mismas hipótesis que~$g$, se reduce la cuestión a demostrar~\Ref{IX.4.9}
para~$S_1$ en lugar de~$S$. Tomando primero como~$A_1$ el producto directo
de los $A/\goth{p}_i$, para los ideales primos mínimos~$\goth{p}_i$
de~$A$, se reduce al caso en que~$A$ es un \emph{dominio de integridad}.
Se demuestra entonces\footnote{\Cf EGA IV~14.3.13 y~14.5.4.} que existe un
subesquema íntegro~$S_1$ de~$S'$, cuasifinito sobre~$S$ y dominante
sobre~$S$, que pasa por un punto de la fibra de~$S'$ sobre el punto
cerrado~$y$ de~$S$ (gracias a que~$S'$ es universalmente abierto y de tipo
finito sobre el esquema local noetheriano íntegro~$S$, y a que
$S_y'\neq\emptyset$). Como~$A$ es completo, $S_1$ es finito sobre~$S$ y,
como domina~$S$, el morfismo $S_1\to S$ es sobreyectivo. Sustituyendo una
vez más~$S$ por~$S_1$, se reduce al caso en que~$S'$ tiene una sección
sobre~$S$, caso en el cual el enunciado es trivial.

\begin{theoreme}
\label{IX.4.10}
Sea $g\colon S'\to S$ un morfismo finito, radicial, sobreyectivo y de
presentación finita (esta última condición es superflua si~$S$ es
localmente noetheriano\footnote{Basta incluso que~$g$ sea entero, radicial
y sobreyectivo, como se ve mediante una reducción fácil al caso del texto,
al estilo de EGA~IV~8; \cf SGA~4 VIII~1.1.}).
\marginpar{241}
Entonces el funtor de imagen inversa induce una equivalencia de la categoría
de los preesquemas étales sobre~$S$ con la categoría de los preesquemas
étales sobre~$S'$.
\end{theoreme}

Como los morfismos diagonales de~$S'$ en $S'\times_S S'$ y
$S'\times_S S'\times_S S'$ son inmersiones sobreyectivas y, por tanto,
inducen en virtud de~\Ref{IX.1.9} equivalencias de las categorías de los
preesquemas étales sobre~$S'\times_S S'$ \resp
$S'\times_S S'\times_S S'$ con la categoría de los preesquemas étales
sobre~$S'$, se deduce de los argumentos formales del descenso~\cite{IX.D}
que todo~$X'$ étale sobre~$S'$ está provisto de un único dato de descenso
respecto de $g\colon S'\to S$. Por tanto,~\Ref{IX.3.3} implica que el
funtor de imagen inversa por~$g$, de la categoría de los preesquemas étales
sobre~$S$ en la categoría de los preesquemas étales sobre~$S'$, es
\emph{plenamente fiel}. Queda por demostrar que es esencialmente
sobreyectivo, es decir, que todo~$X'$ étale sobre~$S'$ es isomorfo a la
imagen inversa de un~$X$ étale sobre~$S$. Como la cuestión es evidentemente
local sobre~$S$ \emph{y sobre}~$X'$, se puede suponer que $S$,~$S'$ y~$X'$
son afines. Pero entonces~$X'$ está separado y es de tipo finito
sobre~$S'$, y se puede aplicar el criterio de efectividad~\Ref{IX.4.7}.

\begin{corollaire}
\label{IX.4.11}
La conclusión de
\ifthenelse{\boolean{orig}}
{~\Ref{IX.4.9}}
{~\Ref{IX.4.10}}
sigue siendo válida si se sustituye la hipótesis sobre~$g$ por la
siguiente: $g$ es fielmente plano, cuasicompacto y radicial.
\end{corollaire}

La misma demostración, invocando~\Ref{IX.4.1} en lugar
de~\Ref{IX.4.7}.

Se observará que la demostración de~\Ref{IX.4.7} es
\og elemental\fg{} en cuanto no utiliza los teoremas
\ifthenelse{\boolean{orig}}
{}
{de}
finitud y de comparación para los morfismos propios (EGA~III~3,~4,~5).
No sucede ya lo mismo con el resultado siguiente:

\begin{theoreme}
\label{IX.4.12}
Sea $g\colon S'\to S$ un morfismo propio, sobreyectivo y de presentación
finita (esta última hipótesis es consecuencia de la primera si~$S$ es
localmente noetheriano). Entonces~$g$ es un morfismo de descenso efectivo
para la categoría fibrada de los recubrimientos étales de preesquemas.
\end{theoreme}

En virtud de~\Ref{IX.3.3} y de~\Ref{IX.2.2}, se reduce la cuestión a
demostrar que, para todo recubrimiento étale~$X'$ sobre~$S'$, provisto de
un dato de descenso respecto de $g\colon S'\to S$, este dato es efectivo.
Utilizando~\Ref{IX.4.3}, se reduce fácilmente
\marginpar{242}
al caso en que~$S$ es noetheriano y, utilizando~\Ref{IX.4.6}, se puede
suponer, por tanto, que~$S$ es el espectro de un anillo local noetheriano
\emph{completo}~$A$. Introduzcamos~$S''$ y~$S'''$ como de costumbre; sea
$S_0$ el espectro del cuerpo residual de~$A$, y sean~$S_0'$,~$S_0''$ y
$S_0'''$ los preesquemas deducidos de~$S'$,~$S''$ y~$S'''$ por el cambio
de base $S_0\to S$, es decir, las fibras de~$S'$,~$S''$ y~$S'''$ en el
punto cerrado de~$S$. Según~\Ref{IX.1.10}, los morfismos $S_0\to S$,
$S_0'\to S'$, etc., inducen equivalencias de la categoría de los
recubrimientos étales sobre el esquema de llegada con la categoría de los
recubrimientos étales sobre el esquema de origen. En consecuencia,
$g\colon S'\to S$ es un morfismo de descenso estricto para la categoría
fibrada de los recubrimientos étales de preesquemas si y solo si lo es
$g_0\colon S_0'\to S_0$, lo que, en efecto, ocurre en virtud
de~\Ref{IX.4.1}. Esto concluye la demostración de~\Ref{IX.4.12}. (En este
razonamiento solo se necesitaba de~\Ref{IX.1.10} el hecho de que el funtor
considerado allí es \emph{plenamente fiel}, lo cual \emph{no} utiliza el
teorema de existencia de haces coherentes en geometría algebraica.)
